% Created 2026-03-09 Mon 22:06
% Intended LaTeX compiler: pdflatex
\documentclass[11pt]{article}
\usepackage[utf8]{inputenc}
\usepackage[T1]{fontenc}
\usepackage{graphicx}
\usepackage{longtable}
\usepackage{wrapfig}
\usepackage{rotating}
\usepackage[normalem]{ulem}
\usepackage{amsmath}
\usepackage{amssymb}
\usepackage{capt-of}
\usepackage{hyperref}
\author{The Prologos Project}
\date{2026-02-24}
\title{Design Methodology\\\medskip
\large From Research to Realization — How We Design and Build Prologos}
\hypersetup{
 pdfauthor={The Prologos Project},
 pdftitle={Design Methodology},
 pdfkeywords={},
 pdfsubject={},
 pdfcreator={Emacs 30.2 (Org mode 9.7.39)}, 
 pdflang={English}}
\begin{document}

\maketitle
\tableofcontents

\section{Purpose}
\label{sec:org38268c4}

This document describes the methodology we follow when designing and
implementing significant features in Prologos. It is both a process guide
and a set of meta-lessons about \emph{how to do design work well} in a language
with deep theoretical foundations and practical engineering constraints.

The methodology emerged organically from our practice on major efforts
including the Extended Spec Language Design, the Logic Engine on
Propagators, Collections Ergonomics, the Numerics Tower, and the
Homoiconicity phases. It is descriptive (capturing what actually works)
as much as prescriptive (defining what we should do).
\section{The Five Stages}
\label{sec:org018db5b}

\subsection{Stage 1: Deep Research}
\label{sec:org29a7a1b}

\emph{Goal}: Ground ourselves in the best techniques --- cutting-edge and
well-proven --- with a comprehensive survey of the landscape.
\subsubsection{What this looks like}
\label{sec:org9715882}

\begin{itemize}
\item Read primary literature (papers, specifications, reference implementations)
\item Survey existing systems that address the same problem space
\item Identify the theoretical foundations that our implementation will rest on
\item Produce a research document (\texttt{docs/research/} or \texttt{docs/tracking/}) that
synthesizes findings
\end{itemize}
\subsubsection{Key practices}
\label{sec:orgfef096f}

\begin{itemize}
\item \textbf{Cast a wide net}: Survey at least 3--5 existing approaches to the same
problem. For the Extended Spec Design, we surveyed Clojure Spec, Malli,
PropEr, Agda, Idris 2, Lean 4, Haskell, F*, Koka, and Racket Contracts
before writing a line of design prose.

\item \textbf{Identify the essential vs. accidental}: What is fundamental to the
problem domain vs. what is incidental to a particular implementation?
The Logic Engine research distilled that propagator networks, lattice
theory, and truth maintenance are the essential components --- the
specific data structures (CHAMP vs. HAMTs, persistent vs. mutable)
are implementation choices.

\item \textbf{Name the theories}: Lattice theory, Curry-Howard correspondence,
the pi-calculus, Radul-Sussman propagators, Kuper's LVars --- naming
the theoretical foundations makes them citable, shareable, and
critiquable. A design grounded in named theory is more robust than
one grounded in intuition alone.

\item \textbf{Document as you go}: Research that isn't written down evaporates
when context windows reset. The research document becomes the
institutional memory.
\end{itemize}
\subsubsection{Artifacts}
\label{sec:orgf46db06}

\begin{itemize}
\item Research document: \texttt{docs/research/TOPIC.md} or
\texttt{docs/tracking/YYYY-MM-DD\_TOPIC\_RESEARCH.md}
\item Bibliography of key papers and systems surveyed
\item A vocabulary of concepts that the team will use in subsequent phases
\end{itemize}
\subsubsection{Example}
\label{sec:orgb3fbe5c}

The Logic Engine began with \texttt{2026-02-23\_LATTICE\_PROPAGATOR\_RESEARCH.md},
covering lattice foundations, LVars/LVish deterministic parallelism, the
"Multiverse Mechanism" and "Pocket Universe" theories, Radul-Sussman
propagator networks, and TMS/ATMS truth maintenance. This document
established the conceptual vocabulary (cells, propagators, quiescence,
monotonic merge) that all subsequent design work referenced.
\subsection{Stage 2: Research Refinement and Gap Analysis}
\label{sec:org286b5a1}

\emph{Goal}: Refine our understanding, identify gaps in current infrastructure,
consider tradeoffs, identify opportunities aligned with core principles,
and make qualified recommendations.
\subsubsection{What this looks like}
\label{sec:org162a3db}

\begin{itemize}
\item Review research findings against our existing implementation
\item Identify what infrastructure we already have and what is missing
\item Enumerate alternative approaches with their tradeoffs
\item Check alignment with design principles in \texttt{DESIGN\_PRINCIPLES.org} and
\texttt{LANGUAGE\_VISION.org}
\item Make recommendations, with rationale, for which approach to pursue
\end{itemize}
\subsubsection{Key practices}
\label{sec:org0cdb1ad}

\begin{itemize}
\item \textbf{Infrastructure gap analysis}: Before designing the Logic Engine, we
identified that we needed persistent data structures (CHAMP), lattice
traits, and type-system exposure of runtime values --- and that we
already had CHAMP maps, the trait system, and the 14-file AST pipeline.
The gap analysis shapes the roadmap.

\item \textbf{Tradeoff matrices}: When multiple approaches exist, make the tradeoffs
explicit. Mutable vs. persistent propagator networks? The tradeoff is
performance (mutable wins for raw speed) vs. backtracking simplicity
(persistent wins --- O(1) backtrack is "keep old reference"). Making
this explicit lets us choose with eyes open.

\item \textbf{Principle alignment check}: Every recommendation should pass the
"does this uphold our principles?" test:
\begin{itemize}
\item Does it maintain homoiconicity?
\item Does it support progressive disclosure?
\item Does it decomplect orthogonal concerns?
\item Does it compose with the existing layered architecture?
\item Does it follow the "most generalizable interface" principle?
\end{itemize}

The Logic Engine's decision to subsume LVars into propagator cells
(rather than having separate LVar types) was driven by
decomplection: cells with per-cell merge functions are the more
general abstraction.

\item \textbf{Opportunities over features}: Don't just fill gaps --- identify
opportunities where our unique combination of features enables
something no other system offers. The fusion of propagators with
session types for distributed constraint solving is an opportunity
that arises from having both systems in the same language.
\end{itemize}
\subsubsection{Artifacts}
\label{sec:org2fa224c}

\begin{itemize}
\item Gap analysis section in the design document
\item Tradeoff tables with explicit criteria
\item Recommendation with rationale tied to principles
\end{itemize}
\subsection{Stage 3: Design Iteration}
\label{sec:org6ce9496}

\emph{Goal}: Produce a concrete design with a detailed, phased roadmap for
implementation, refined through critique and feedback until there is
full clarity.
\subsubsection{What this looks like}
\label{sec:org48d4af2}

This is an iterative cycle:

\begin{enumerate}
\item \textbf{Draft}: Produce a complete design document with implementation roadmap
\item \textbf{Critique}: Subject the design to rigorous, independent critique
\item \textbf{Respond}: Address critiques --- accept, refine, or justify
\item \textbf{Repeat}: Continue until all parties have clarity and confidence
\end{enumerate}
\subsubsection{Key practices}
\label{sec:orgb15116b}

\begin{itemize}
\item \textbf{Comprehensive first drafts}: The first draft should be as complete
as possible. Include data structures, type signatures, AST nodes,
phase dependencies, test strategies, and concrete code examples.
An incomplete draft invites incomplete critique.

\item \textbf{Invite adversarial critique}: Seek critique that probes weaknesses,
not confirmation that validates strengths. The Logic Engine design
was subjected to a 10-point critique covering resolution strategy
(underspecified), contradiction detection (hand-wavy), negation
handling (unsound without stratification), and performance analysis
(missing). Every point led to a design improvement.

\item \textbf{Distinguish design issues from implementation issues}: A critique
that says "the resolution strategy is unclear" is a design issue
that must be resolved before implementation. A critique that says
"CHAMP lookup is O(log₃₂ n)" is an implementation detail that can
be optimized later. Don't let implementation concerns block design
progress, but don't let design ambiguity persist into implementation.

\item \textbf{Pushback with context}: Not all critiques are valid. When a critique
misunderstands the system's context, push back with the fuller picture.
When the Logic Engine critique suggested adding \texttt{:includes} to traits
for contradiction detection, we pushed back: "We don't want trait
hierarchies AT ALL, EVER! This is the whole rationale for the \texttt{bundle}
concept." The pushback refined both the design and our articulation of
the principle.

\item \textbf{Phase dependencies are architecture}: The phased roadmap isn't just
a schedule --- it's an architectural statement about what depends on
what. The Logic Engine's critical path
(\texttt{Phase 1 → 2 → 3 → 5 → 6 → 7}, with \texttt{Phase 4} parallelizable)
encodes fundamental architectural dependencies. Getting this right
prevents wasted work.

\item \textbf{Concrete over abstract}: Show elaboration examples, not just type
signatures. Show how \texttt{defr ancestor} becomes propagator cells and
fire functions, not just "relations elaborate to propagator networks."
Every level of abstraction should have at least one concrete example
that a reader can trace through.
\end{itemize}
\subsubsection{The critique cycle in practice}
\label{sec:org2c0d6d3}

The Extended Spec Design went through multiple rounds:

\begin{enumerate}
\item \textbf{Initial research survey} → comprehensive landscape analysis
\item \textbf{First design draft} → keyword symmetry table, metadata map design,
\texttt{property} as composable proposition groups
\item \textbf{Feedback round} → Should \texttt{property} use \texttt{:includes} or \texttt{:extends}?
How do properties attach to traits? Resolution: \texttt{:includes} for
conjunctive composition (set union), \texttt{:laws} on trait for attachment.
\item \textbf{Final refinement} → phased roadmap separating syntax (Phase 1) from
backends (Phase 2+), acknowledging what can be built now vs. what
needs future infrastructure.
\end{enumerate}

The Logic Engine Design went through a similar cycle:

\begin{enumerate}
\item \textbf{Research} → lattice/propagator theory survey
\item \textbf{Vision discussion} → relational language syntax, solver architecture,
provenance, tabling-by-default
\item \textbf{First design draft} → 7-phase roadmap, data structures, AST nodes
\item \textbf{Independent critique} → 10 points of feedback, including fundamental
gaps in resolution strategy and functional-relational boundary
\item \textbf{Response and refinement} → accepted \texttt{HasTop} trait (Point 3), rejected
trait hierarchy (Point 3 alternative), clarified UnionFind vs. cells,
added resolution strategy section, added performance expectations
\item \textbf{Final roadmap} → phased implementation plan with clear phase
dependencies and test strategies per phase
\end{enumerate}
\subsubsection{Artifacts}
\label{sec:org0669fd2}

\begin{itemize}
\item Design document: \texttt{docs/tracking/YYYY-MM-DD\_TOPIC\_DESIGN.org}
(org-mode preferred for structured prose with code blocks)
\item Phased implementation roadmap with explicit dependencies
\item Phase-specific test strategies and success criteria
\item Record of critique and responses (preserved in standups)
\end{itemize}
\subsection{Stage 4: Implementation}
\label{sec:orgaddc453}

\emph{Goal}: Execute the phased roadmap, shipping complete and sound solutions
at each phase, raising design issues as they surface.
\subsubsection{What this looks like}
\label{sec:orga3208de}

\begin{itemize}
\item Work through the roadmap phase by phase
\item Create tracking documents before implementation begins
\item Ship complete, tested solutions --- not partial scaffolding
\item When implementation reveals design gaps, stop and address them
\end{itemize}
\subsubsection{Key practices}
\label{sec:org989217c}

\begin{itemize}
\item \textbf{Completeness over deferral}: This is our most important implementation
principle (see \texttt{DEVELOPMENT\_LESSONS.org} § "Completeness Over Deferral").
When you have clarity, vision, and full context --- finish the work now.
Half-built pieces that get deferred are half-built pieces that get
forgotten. The cost of re-acquiring context later almost always exceeds
the cost of doing the work while the understanding is fresh.

Defer \emph{only} when there is a genuine dependency on unbuilt infrastructure
or genuinely uncertain design. "We'll come back to it" is a red flag.

\item \textbf{Design issues surface during implementation}: This is expected and
healthy. The Logic Engine Phase 3 revealed that:
\begin{itemize}
\item Parametric \texttt{impl} dispatch failed for compound type args without
\texttt{where} constraints (a \texttt{macros.rkt} fix was needed)
\item Generic trait accessor calls inside closures can't resolve implicit
type params (a genuine meta-resolution limitation, deferred with
explicit tracking)
\item QTT multiplicity violations arise from inline lambdas that reference
prelude functions differently from named definitions
\end{itemize}

Each of these was addressed: the first two with code fixes and design
documentation, the third with a workaround and a tracked limitation.
The key is: \emph{don't paper over design issues with quick fixes that meet
criteria partially}.

\item \textbf{If a gap requires deeper infrastructure, address the core concern}:
When the Logic Engine critique identified that contradiction detection
needed a \texttt{HasTop} trait, we built it --- including a \texttt{BoundedLattice}
bundle and trait instances for Bool, FlatVal, and Interval. This was
more work than a simple boolean flag, but it produced the correct
abstraction that composes with the rest of the trait system.

\item \textbf{Phase-gated implementation with sub-phases}: Break phases into
lettered sub-phases (a, b, c\ldots{}) with explicit "done" vs.
"remaining" tracking. The Logic Engine Phase 3 was broken into:
\begin{itemize}
\item 3a: AST structs + mechanical traversals
\item 3b: Type rules + type tests
\item 3c: Reduction rules + eval tests
\item 3d: Surface syntax + integration tests
\item 3e: Library wrappers + trait tests
\end{itemize}

Each sub-phase has a clear scope, produces testable artifacts, and
creates a natural commit point.

\item \textbf{Test at every boundary}: Every sub-phase should end with passing
tests. The Logic Engine implementation produced 56 new tests across
3 test files, verifying types (32 tests), integration (16 tests),
and library-level traits (8 tests). Tests are not an afterthought
--- they are the proof that the implementation matches the design.

\item \textbf{Track deferred work immediately}: When something must be deferred
(the \texttt{new-lattice-cell} generic wrapper, for instance), add it to
\texttt{DEFERRED.md} in the same commit. Deferred work not tracked is
abandoned work.
\end{itemize}
\subsubsection{Implementation flow}
\label{sec:org27d9960}

\begin{verbatim}
┌──────────────────────────────────┐
│ Create tracking document         │
│ (before writing any code)        │
└───────────────┬──────────────────┘
                │
┌───────────────▼──────────────────┐
│ Sub-phase a: Foundation          │◄──── Tests pass? ──── Yes ──► Commit
│ (structs, types, basic cases)    │         │
└───────────────┬──────────────────┘         No
                │                            │
┌───────────────▼──────────────────┐         ▼
│ Sub-phase b: Core logic          │    Fix / Investigate
│ (type rules, reduction, etc.)    │         │
└───────────────┬──────────────────┘         │
                │                  ◄─────────┘
                ▼
┌──────────────────────────────────┐     ┌──────────────────┐
│ Design issue discovered?         │────►│ Stop. Address it. │
│                                  │ Yes │ Update design doc │
└───────────────┬──────────────────┘     │ Then continue.    │
                │ No                     └──────────────────┘
                ▼
┌──────────────────────────────────┐
│ Sub-phases c, d, e...            │
│ (repeat until phase complete)    │
└───────────────┬──────────────────┘
                │
┌───────────────▼──────────────────┐
│ Full test suite + whale check    │
│ Update DEFERRED.md               │
│ Update tracking document         │
│ Commit                           │
└──────────────────────────────────┘
\end{verbatim}
\subsubsection{Artifacts}
\label{sec:org362ac80}

\begin{itemize}
\item Phase-specific commits with descriptive messages
\item Updated tracking document with completion status
\item Updated \texttt{DEFERRED.md} for any deferred work
\item Test suite growth (measurable)
\end{itemize}
\subsection{Stage 5: Composition and Extension}
\label{sec:org8d3205f}

\emph{Goal}: Enjoy the well-thought-out design and how it composes with our
multi-level, modular, extensible language. Verify that the new feature
integrates cleanly and enables future extension.
\subsubsection{What this looks like}
\label{sec:org2c45c76}

\begin{itemize}
\item Write a capstone demo (\texttt{.prologos} file) exercising the feature in
the user-facing syntax — this is the first mandatory step
\item Verify that the new feature composes with existing features
\item Check that the new abstractions are reusable in contexts beyond the
original design
\item Update principles documents if new patterns emerged
\item Identify the new possibilities that the feature enables
\item Conduct a Post Implementation Review for significant features
\end{itemize}
\subsubsection{Key practices}
\label{sec:orgcfb9ec9}

\begin{itemize}
\item \textbf{Composition is the test}: A well-designed feature should compose
with features it wasn't specifically designed for. The propagator
network composes with the trait system (lattice traits define merge
behavior), with the type system (PropNetwork is a first-class type),
and with the prelude (users can write lattice-aware programs with
standard library functions). If a feature \emph{doesn't} compose, that's
a design smell worth investigating.

\item \textbf{Extension paths are visible}: After shipping Phase 3 of the Logic
Engine, the path to Phase 4 (UnionFind), Phase 5 (ATMS), and
Phase 7 (surface syntax) is clear. Each depends on the infrastructure
just built, and the design document already specifies what each will
need. A good implementation \emph{opens doors} rather than closing them.

\item \textbf{Update the living documents}: When a feature completes, update:
\begin{itemize}
\item \texttt{DEFERRED.md} (mark phases complete, add new deferrals)
\item \texttt{MEMORY.md} (project memory for context recovery)
\item Relevant principles docs (new patterns, new lessons)
\item The original tracking document (completion status, lessons learned)
\end{itemize}

\item \textbf{Teach through the capstone demo}: Write a \texttt{.prologos} demo file that
shows the feature in context (see "Feature Capstone Demo" below). The
demo serves triple duty: user-facing documentation, integration
validation, and honest assessment of WS-mode surface coverage. This
replaces the test-files-as-documentation pattern — tests prove
correctness, demos prove usability.
\end{itemize}
\subsubsection{Feature Capstone Demo (Mandatory)}
\label{sec:orgb267be6}

After a significant feature track is complete --- all sub-phases shipped,
tests passing --- write a \textbf{capstone demo file} that exercises the feature
in the user-facing \texttt{.prologos} syntax. This is not a test file (those
exist already); it is a \emph{feature showcase} written in the ideal syntax of
the language, intended to be read and run by users (including future selves).
\begin{enumerate}
\item Why this is mandatory
\label{sec:org758e5fb}

The capstone demo serves purposes that unit tests cannot:

\begin{itemize}
\item \textbf{Discovers syntax gaps}: Unit tests exercise the sexp/API layer. The
capstone forces the WS reader → preparse → parse → elaborate path,
revealing features that work internally but aren't surfaced in the
user-facing syntax. The FL Narrowing capstone revealed that ground
constructor args in WS narrowing (\texttt{[add zero ?y] = 5N}), nested match
patterns, and composed-function narrowing don't yet work in WS mode ---
none of which were visible from the passing test suite.

\item \textbf{Validates integration with other features}: The demo should exercise
how the new feature interacts with existing language components ---
the prelude, the type system, other syntactic forms (pipes, match,
pattern defn). A feature that works in isolation but breaks when
composed with the rest of the language is not done.

\item \textbf{Establishes ground truth for "what works"}: After the demo runs,
you have an honest inventory: which ideal-syntax expressions produce
correct output, which need commenting out, and which reveal future
work. This inventory is more valuable than any design document for
understanding the current state of the feature.

\item \textbf{Creates user-facing documentation}: The demo is a teaching artifact.
It shows users what the feature does, how to use it, and what the
output looks like. Heavy commenting explains the underlying mechanisms.

\item \textbf{Exercises the language vision}: Writing the demo forces you to think
about how the feature \emph{should} look in the ideal syntax, even if the
current parser doesn't support it yet. Ideal syntax that's commented
out with a \texttt{NOTE:} becomes a specification for future work.
\end{itemize}
\item What the capstone should contain
\label{sec:orgf30f295}

\begin{itemize}
\item \textbf{Prologos header and \texttt{ns} declaration}: First-class \texttt{.prologos} file
in the examples directory
\item \textbf{Sections organized by feature/phase}: Each section introduces one
capability with commentary explaining the mechanism
\item \textbf{Live executable expressions}: Every uncommented expression must run
without errors
\item \textbf{Commented ideal syntax}: Features not yet surfaced in WS mode are
shown as commented-out ideal syntax with \texttt{NOTE:} explaining the gap
\item \textbf{Cross-feature exercises}: At least one section showing how the new
feature composes with existing features (e.g., narrowing + functional
code, narrowing + prelude functions, narrowing + user definitions)
\item \textbf{Summary table}: Lines/tests per sub-phase, giving readers a sense
of the implementation's scope
\end{itemize}
\item What the capstone should NOT be
\label{sec:org2e79e47}

\begin{itemize}
\item Not a test file (no \texttt{check-equal?}, no \texttt{rackunit})
\item Not an exhaustive regression suite (that's what \texttt{tests/} is for)
\item Not a design document (that's in \texttt{docs/tracking/})
\item Not a copy of existing tests translated to WS mode (it should explore
NEW combinations and usage patterns)
\end{itemize}
\item When to write it
\label{sec:org0f316d9}

After the last sub-phase is committed and all tests pass, but \emph{before}
the PIR. The capstone often surfaces issues that inform the PIR's
"challenges" and "forward enablement" sections.
\item Naming and location
\label{sec:org81f05de}

\texttt{racket/prologos/examples/\{feature\}-demo.prologos}

Examples: \texttt{narrowing-demo.prologos}, \texttt{relational-demo.prologos},
\texttt{session-demo.prologos}.
\item Example
\label{sec:org78a4ea7}

The FL Narrowing capstone (\texttt{examples/narrowing-demo.prologos}, 533 lines)
covered all 11 sub-phases with 16 live narrowing queries and 5 commented
ideal-syntax sections. It revealed 3 WS-mode gaps, validated 7 working
cross-feature interactions, and established the ground truth that
narrowing currently operates over inductively-defined types (Nat, Bool)
with Int narrowing via constraint propagation as a future extension.
\end{enumerate}
\subsubsection{Post Implementation Review (PIR)}
\label{sec:org5e44967}

After a significant feature is complete --- all sub-phases shipped, tests
passing, tracking documents updated --- conduct a Post Implementation
Review. The PIR is the reflective counterpart to the research stage: where
Stage 1 asks "what do we know going in?", the PIR asks "what do we know
coming out?"

A PIR is \emph{not} a retrospective about process (though process insights
belong). It is primarily a \emph{technical} review: did the design hold up?
What did the implementation reveal about our architecture? What does
this feature enable that we didn't anticipate?
\begin{enumerate}
\item What to cover
\label{sec:orgc0cd1e1}

\begin{itemize}
\item \textbf{Objectives vs. actuals}: What did we set out to build? What did we
actually build? Quantify scope adherence (e.g., "36/38 sub-phases,
95\%"). Note what was correctly deferred vs. what was scope creep or
missed.

\item \textbf{What went well}: Identify decisions, patterns, and practices that
paid off. These are candidates for codification in principles
documents. The Session Type PIR identified the 3-stage design process
(research → design → plan) and propagator-as-universal-substrate as
key successes worth replicating.

\item \textbf{Challenges and how they were resolved}: Document the hard problems
and their solutions. Be specific --- "the \texttt{??} / \texttt{\$typed-hole} reader
conflict required context-sensitive disambiguation at three points in
the desugaring pipeline" is more useful than "WS parsing was tricky."
Future implementations facing similar challenges will benefit from
the concrete resolution.

\item \textbf{Architectural validation}: Did the feature integrate cleanly with
existing systems? Were the extension points (AST pipeline, trait
system, propagator network) sufficient, or did they need modification?
A clean integration validates the architecture; friction points
identify areas for improvement.

\item \textbf{Quantitative summary}: Lines of code, test count growth, number of
commits, implementation duration. These establish baselines for
estimating future work of similar scope.

\item \textbf{Forward enablement}: What does this feature make possible that wasn't
possible before? What are the natural next steps? This section bridges
from the completed work to the future roadmap.

\item \textbf{Lessons learned}: Distill the key insights --- both technical and
methodological --- into actionable guidance. The best lessons are
specific enough to apply ("single-method trait dicts should be the
function itself, not a wrapper struct") rather than generic ("test
early and often").
\end{itemize}
\item When to do a PIR
\label{sec:org97782ff}

Not every change warrants a PIR. Use one when:
\begin{itemize}
\item The feature spanned multiple phases or sub-phases
\item The implementation took more than a single session
\item The work involved architectural decisions (new AST nodes, new
infrastructure files, new cross-cutting concerns)
\item You learned something that future work should benefit from
\end{itemize}

Skip for routine bug fixes, single-file changes, or features that
followed established patterns without surprises.
\item Artifacts
\label{sec:orge21090e}

\begin{itemize}
\item PIR document: \texttt{docs/tracking/YYYY-MM-DD\_TOPIC\_PIR.md}
\item Updates to \texttt{DEVELOPMENT\_LESSONS.org} for broadly applicable lessons
\item Updates to \texttt{PATTERNS\_AND\_CONVENTIONS.org} for new patterns discovered
\item Updates to \texttt{MEMORY.md} for project-level status
\end{itemize}
\item Example
\label{sec:orgf488a02}

The Session Types PIR (\texttt{2026-03-04\_SESSION\_TYPE\_PIR.md}) covered 36/38
sub-phases across \textasciitilde{}24 hours. Key findings:
\begin{itemize}
\item The 3-stage design methodology (deep research → adversarial critique →
phased implementation) produced zero architectural rework
\item Propagator networks proved to be a universal substrate: session types
were the third domain (after types and multiplicities) to use the
same infrastructure without modification
\item The \texttt{??} / \texttt{\$typed-hole} WS reader conflict revealed that
context-sensitive token disambiguation requires handling at every
stage of the desugaring pipeline, not just one
\item Phase 0 scoping (defer concurrent runtime) was validated as correct:
complete type-level infrastructure enables future runtime work without
redesign
\end{itemize}
\end{enumerate}
\section{Cross-Cutting Principles}
\label{sec:org92c4466}

These principles apply across all five stages.
\subsection{Theoretical Grounding, Practical Surface}
\label{sec:org0858d94}

Every design decision should be traceable to a theoretical foundation,
but the surface presentation should be practical and approachable.
"Lattice-theoretic monotonic constraint propagation" is the theory;
\texttt{net-cell-write} with a merge function is the practice. Both must
exist: the theory ensures correctness, the practice ensures usability.
\subsection{Design Documents Are Living}
\label{sec:org639c8d5}

Design documents are not write-once artifacts. They are updated as
implementation reveals new information. The Logic Engine design document
was refined after every critique round and after Phase 3 implementation
revealed the \texttt{HasTop} trait need and the implicit-resolution limitation.
Stale design documents are worse than no design documents.
\subsection{Critique Is a Gift}
\label{sec:org8759102}

Adversarial critique --- even harsh critique --- improves designs. The
10-point critique of the Logic Engine identified real gaps (resolution
strategy underspecified, negation handling unsound). Every accepted
critique made the design stronger. The discipline is to engage with
critique on its merits, not defensively.

Not all critiques are valid. The key judgment is: does this critique
identify a real gap in the design, or does it misunderstand the
system's context? When the latter, push back clearly and use the
pushback to sharpen your articulation of \emph{why} the design is the way
it is.
\subsection{Standups as Design Memory}
\label{sec:orgd4ce6b5}

Our standup documents (\texttt{docs/standups/standup-YYYY-MM-DD.org}) serve as
a chronological record of design discussions, including the back-and-forth
of critique and refinement. The 🗣️ (human) and 🤖 (machine) annotations
preserve the conversational flow. This record is invaluable for
reconstructing design rationale months later.
\subsection{The 14-File Pipeline as Discipline}
\label{sec:orgaca980a}

Prologos's 14-file AST pipeline (syntax → surface-syntax → parser →
elaborator → typing-core → qtt → reduction → substitution → zonk →
pretty-print → unify → macros → foreign) is both a cost and a benefit.
The cost is that every new AST node requires touching many files. The
benefit is that every subsystem handles every node consistently. This
enforced consistency is a design methodology in itself: if a feature
can't be expressed through the pipeline, it's a signal that the
feature's abstraction is wrong.
\section{Anti-Patterns}
\label{sec:org7a8f38d}

\subsection{"We'll come back to it"}
\label{sec:org2ed5d93}

Defer only when genuinely blocked. Track all deferrals immediately.
See \texttt{DEVELOPMENT\_LESSONS.org} § "Completeness Over Deferral."
\subsection{Quick fixes that meet criteria partially}
\label{sec:orgea3405b}

A partial solution that passes tests but doesn't address the underlying
design concern creates technical debt with interest. Better to identify
the design gap, address it, and ship the complete solution.
\subsection{Research without documentation}
\label{sec:org65dd746}

Research that exists only in a developer's head is lost when context
resets. Write it down. Even rough notes are better than nothing.
\subsection{Design without critique}
\label{sec:org26e7da9}

A design that hasn't been subjected to adversarial critique has
unknown weaknesses. Actively seek feedback, especially on the parts
you're most confident about --- confidence is where blind spots hide.
\subsection{Implementation without tests}
\label{sec:org45dab49}

An untested implementation is an unverified hypothesis. Tests are the
proof that the implementation matches the design. New features aren't
done until they have tests.
\section{Relationship to Other Principles Documents}
\label{sec:org0898d83}

\begin{center}
\begin{tabular}{ll}
Document & Relationship to this methodology\\
\hline
\texttt{DESIGN\_PRINCIPLES.org} & Provides the values that Stage 2 alignment checks use\\
\texttt{LANGUAGE\_VISION.org} & Provides the north star that research (Stage 1) aims at\\
\texttt{DEVELOPMENT\_LESSONS.org} & Collects lessons that emerge from Stage 4\\
\texttt{PATTERNS\_AND\_CONVENTIONS.org} & Captures patterns that Stage 5 discovers\\
\texttt{ERGONOMICS.org} & Informs Stage 2 tradeoff analysis for user-facing features\\
\texttt{RELATIONAL\_LANGUAGE\_VISION.org} & Example of Stage 1-2 output for a specific domain\\
\end{tabular}
\end{center}
\section{Exemplar Projects}
\label{sec:orgd078957}

These completed efforts illustrate the methodology in practice:
\subsection{Extended Spec Language Design}
\label{sec:orgdff6ff9}
\begin{itemize}
\item \textbf{Research}: Survey of 10+ systems (Clojure Spec, Malli, PropEr, Agda,
Idris 2, Lean 4, Haskell, F*, Koka, Racket Contracts)
\item \textbf{Design}: \texttt{2026-02-22\_EXTENDED\_SPEC\_DESIGN.org} with keyword symmetry
table, metadata map design, 4-phase roadmap
\item \textbf{Iteration}: Multiple rounds on \texttt{property} composition, \texttt{:includes}
semantics, functor metadata
\item \textbf{Implementation}: Phase 1 complete (\texttt{??} holes, \texttt{property}, \texttt{functor},
\texttt{:examples}, \texttt{:deprecated}). Phases 2+ tracked in DEFERRED.md.
\item \textbf{Composition}: \texttt{property} composes with \texttt{trait} via \texttt{:laws}, \texttt{bundle}
composes \texttt{property} groups, \texttt{functor} composes with \texttt{spec} metadata.
\end{itemize}
\subsection{Logic Engine on Propagators}
\label{sec:orge8fe517}
\begin{itemize}
\item \textbf{Research}: \texttt{2026-02-23\_LATTICE\_PROPAGATOR\_RESEARCH.md} + the broader
\texttt{TOWARDS\_A\_GENERAL\_LOGIC\_ENGINE\_ON\_PROPAGATORS.org} synthesis
\item \textbf{Design}: \texttt{2026-02-24\_LOGIC\_ENGINE\_DESIGN.org} with 7-phase roadmap,
14 AST nodes, data structure specifications
\item \textbf{Iteration}: Relational language vision discussion, design critique
(10 points), pushback on trait hierarchies, refinement of resolution
strategy and HasTop design
\item \textbf{Implementation}: Phases 1--3 complete. Phase 1 (lattice traits),
Phase 2 (persistent PropNetwork with BSP), Phase 3 (14 AST nodes
through 12-file pipeline, HasTop trait, BoundedLattice bundle).
56 new tests, parametric impl dispatch fix.
\item \textbf{Composition}: PropNetwork composes with trait system (Lattice, HasTop,
BoundedLattice), type system (first-class type), prelude (available
from any \texttt{ns} module).
\end{itemize}
\subsection{Collections Ergonomics}
\label{sec:org02d89a4}
\begin{itemize}
\item \textbf{Research}: Audit of existing collection generics, gap identification
\item \textbf{Design}: Stages A--I with clear boundaries
\item \textbf{Implementation}: Stages A--H complete in one session (completeness
over deferral). Stage I deferred (genuine infrastructure dependency
on transient type exposure).
\item \textbf{Composition}: Generic \texttt{map\textasciitilde{}/\textasciitilde{}filter\textasciitilde{}/\textasciitilde{}reduce} work across all
collection types. \texttt{into} converts between any two collections.
\end{itemize}
\subsection{Session Types (S1--S8)}
\label{sec:org1b8f046}
\begin{itemize}
\item \textbf{Research}: \texttt{2026-03-03\_SESSION\_TYPE\_RESEARCH.md} surveying the
pi-calculus, Gay-Vasconcelos session types, Toninho-Caires-Pfenning
propositions-as-sessions, Wadler's Classical Processes, Scalas-Yoshida
multiparty session types, and Jespersen-Munksgaard-Sestoft Rust
embedding
\item \textbf{Design}: \texttt{2026-03-03\_SESSION\_TYPE\_DESIGN.md} with 8-phase roadmap,
\textasciitilde{}26 surface syntax structs, propagator-based type checking,
Galois connection bridges
\item \textbf{Iteration}: Adversarial critique refined duality computation,
deadlock detection strategy, capability integration, and the
functional-relational boundary between session and process layers
\item \textbf{Implementation}: 36/38 sub-phases complete (\textasciitilde{}24 hours). 4 new
infrastructure files, 1,880 new lines, 392 test cases (+21\% suite
growth). S8b correctly deferred (promises meaningless without
concurrent runtime).
\item \textbf{Composition}: Session types are the third domain on the propagator
network (after types and multiplicities), using the same
infrastructure without modification. Galois connection bridges
compose session and type domains. \texttt{defproc} composes with trait
system, QTT, and capability binders.
\item \textbf{PIR}: \texttt{2026-03-04\_SESSION\_TYPE\_PIR.md} --- first exemplar of the
PIR practice. Validated the 3-stage design methodology, identified
the \texttt{??\textasciitilde{}/\textasciitilde{}\$typed-hole} desugaring lesson, and established forward
roadmap for concurrent runtime and multi-party sessions.
\end{itemize}
\subsection{Effect Ordering via Session Types (Architecture A+D)}
\label{sec:org0401db1}
\begin{itemize}
\item \textbf{Research}: \texttt{2026-03-06\_EFFECTFUL\_PROPAGATORS\_RESEARCH.md} surveying
CALM theorem, LVars, Timely Dataflow, CRDTs, Radul-Sussman, BloomL,
algebraic effects, and Flix. Three candidate architectures (A, B, C)
for recovering order-dependent effects on a monotone substrate.
\item \texttt{2026-03-06\_SESSION\_TYPES\_AS\_EFFECT\_ORDERING.org} (\textasciitilde{}1,200 lines):
discovery that session types are \emph{causal clocks} --- each position in
a session type's continuation chain is a causally-ordered event.
Formal Galois connection \((\alpha, \gamma)\) between session lattice
and effect position lattice. Quantale morphism at the algebraic level.
\item \textbf{Key insight}: The Layered Recovery Principle --- non-monotone behavior
(effect ordering) recovered on a monotone substrate (propagator network)
by inserting a control layer between phases of monotone computation.
Generalises the logic engine's recovery of choice-point semantics.
\item \textbf{Design}: \texttt{2026-03-07\_ARCHITECTURE\_AD\_IMPLEMENTATION\_DESIGN.org}
(\textasciitilde{}750 lines) with 16 sub-phases across 6 phases (AD-A through AD-F),
dependency graph, concrete data structures, traced example, migration
strategy, and test plan (\textasciitilde{}178 tests).
\item \textbf{Architecture decision}: A + D --- session-derived ordering (D) for
session-typed IO, sequential walk ordering (A) for unsessioned IO.
Architecture B subsumed; no overlap, each mechanism maximally suited
to its domain.
\item \textbf{Implementation}: All 6 phases complete. 4 new modules
(\texttt{effect-position.rkt}, \texttt{effect-bridge.rkt}, \texttt{effect-ordering.rkt},
\texttt{effect-executor.rkt}), \texttt{770 lines, \textasciitilde{}165 tests, 0 regressions.
  Auto-detection in \textasciitilde{}driver.rkt} selects Architecture A or D based on
process structure.
\item \textbf{Principles}: Elevated to \texttt{EFFECTFUL\_COMPUTATION\_ON\_PROPAGATORS.org}
--- the Layered Recovery Principle as a general design methodology,
the \texttt{rel\textasciitilde{}/\textasciitilde{}proc} structural parallel, and the five-layer architecture.
\item \textbf{Composition}: Effect position lattice composes with session lattice
via Galois connection bridge propagators. Effect ordering composes
with ATMS branching (choice branches → per-branch orderings).
Architecture selection composes with existing \texttt{rt-execute-process}
(additive, not replacing).
\end{itemize}
\end{document}
